\ChapterImageStar[cap:conclusiones]{Conclusiones}{./images/fondo.png}\label{cap:conclusiones}
\mbox{}\\

El desarrollo de esta investigación permitió alcanzar el objetivo general de proponer una Notación de Diseño de Infraestructura de TI (\NDITI) adaptada al contexto del programa de Ingeniería de Sistemas y Computación de la Universidad del Quindío. A través de un enfoque metodológico estructurado que integró caracterización del contexto académico, estudio de mapeo sistemático (SMS), análisis de decisiones (DAR) y un proceso de validación con estudiantes, se logró diseñar, fundamentar y evaluar una propuesta de notación orientada a mejorar la representación de infraestructuras tecnológicas.

\section{Sobre la metodología empleada}
\noindent

La metodología adoptada en el proyecto combinó múltiples enfoques complementarios que permitieron abordar el problema desde una perspectiva integral. En primer lugar, la caracterización del contexto ---incluyendo docentes, estudiantes y dinámicas del grupo \GRID--- permitió identificar necesidades y problemáticas reales en torno al uso de notaciones para modelado de infraestructura.

Posteriormente, se realizó un estudio de mapeo sistemático (SMS) que permitió identificar referentes metodológicos y tecnológicos relevantes, asegurando que la propuesta se fundamentara en evidencia académica y prácticas existentes. Este proceso proporcionó una base sólida para la selección de notaciones candidatas.

Finalmente, la aplicación de la metodología DAR permitió evaluar de manera estructurada dichas alternativas, considerando criterios como popularidad, documentación, formalidad, vigencia y minimalismo. Este proceso garantizó trazabilidad y rigor en la toma de decisiones.

En conjunto, la integración de estas metodologías demostró ser adecuada para el diseño de propuestas en contextos académicos, al garantizar que cada decisión de diseño tuviera respaldo tanto en la literatura como en las necesidades reales del entorno.

\section{Sobre la construcción de la notación~\NDITI}
\noindent

El proceso de diseño de la~\NDITI~se fundamentó principalmente en la especificación de ArchiMate, seleccionada como base debido a su madurez, estructura multicapa y capacidad para representar arquitecturas de forma integral.

Sin embargo, el análisis realizado evidenció limitaciones en ArchiMate para contextos específicos de infraestructura técnica, así como dificultades de aprendizaje en entornos académicos. Por esta razón, la propuesta~\NDITI~no se limitó a adoptar este lenguaje, sino que incorporó elementos conceptuales de otras notaciones como UML y BPMN, con el objetivo de complementar sus capacidades.

Este enfoque permitió construir una notación híbrida que mantiene la estructura por capas coherente heredada de ArchiMate, mejora la precisión técnica en la representación por influencia de UML, y facilita la comprensión y comunicación de los modelos a partir de principios tomados de BPMN. Como resultado, la~\NDITI~constituye una propuesta adaptada al contexto educativo que busca equilibrar formalidad, expresividad y facilidad de uso.


\section{Sobre la validación de la notación}
\noindent
La validación de la~\NDITI~se realizó mediante un ejercicio aplicado con estudiantes del programa de Ingeniería de Sistemas y Computación, quienes utilizaron la notación en una actividad de diagramación y posteriormente respondieron una encuesta de percepción.

El análisis comparativo entre el grupo que utilizó la notación y el grupo de diagramación libre permitió identificar diferencias relevantes en los cuatro criterios evaluados. En estructuración del modelo, el grupo~\NDITI~obtuvo un promedio de 4.6 frente a 3.0 del grupo libre, constituyendo la diferencia más amplia entre grupos (1.6 puntos) y el criterio con mayor homogeneidad interna en el grupo experimental, donde cuatro de cinco subgrupos alcanzaron la valoración máxima. En correspondencia con el enunciado y completitud del modelo, ambos grupos registraron promedios de 4.6 y 3.6, y 4.6 y 3.4 respectivamente, evidenciando que la notación favorece la fidelidad y exhaustividad de los modelos construidos. En representación de relaciones, la diferencia fue de 4.0 frente a 3.4, siendo este el criterio de mayor variabilidad interna en ambos grupos, lo que sugiere que la semántica de las relaciones representa el principal reto de apropiación tanto con notación como sin ella.

Más allá de los promedios, la variabilidad interna de cada grupo ofrece información diferenciadora: el grupo~\NDITI~presentó un rango de desempeño individual entre 3.75 y 5.0, mientras que el grupo libre osciló entre 2.50 y 4.25. Esta menor dispersión en el grupo experimental indica que la notación no solo eleva el desempeño promedio, sino que establece un piso mínimo de calidad más consistente, independientemente de las capacidades o conocimientos previos de los participantes.

En cuanto a la encuesta de percepción, los resultados evidencian una valoración positiva en los cuatro ejes evaluados. En comprensión de la notación (Eje I), los promedios se ubicaron entre 4.0 y 5.0, destacándose la utilidad de los ejemplos (≈5.0) y la organización por capas (≈4.7) como los aspectos mejor valorados. En uso de los elementos (Eje II), los promedios oscilaron entre 4.0 y 4.25, con una percepción favorable sobre la cantidad y manejabilidad de los elementos disponibles. En expresividad y utilidad (Eje III), este eje concentró las valoraciones más altas de la encuesta, con promedios entre 4.5 y 4.75, siendo la utilidad para modelar sistemas reales el ítem mejor valorado (≈4.75). En experiencia general (Eje IV), los promedios se situaron entre 3.5 y 4.25, con la facilidad de aprendizaje como el ítem de menor valoración (≈3.5), lo que señala la curva de aprendizaje inicial como el principal desafío percibido.

La integración de ambas dimensiones revela una coherencia general entre percepción y desempeño: los aspectos mejor valorados en la encuesta —organización por capas y expresividad— corresponden precisamente a los criterios en que el grupo~\NDITI~mostró mayor ventaja sobre el grupo libre. La única tensión identificada se sitúa en la facilidad de aprendizaje, cuya valoración más baja contrasta con un desempeño consistentemente superior al del grupo de referencia, sugiriendo que la estructura de la notación compensa parcialmente las dificultades iniciales de familiarización.

\section{Sobre el impacto en el contexto académico}
\noindent
La propuesta de la~\NDITI~tiene implicaciones directas en el contexto del programa de Ingeniería de Sistemas y Computación y del grupo \GRID, al proporcionar una herramienta conceptual y práctica para la enseñanza y comprensión de infraestructuras de TI.

Los resultados de la validación permiten establecer una comparación implícita entre el modelado guiado por la notación y el modelado sin marco de referencia definido. El grupo libre, utilizado como proxy del estado previo a la adopción de una notación estructurada, evidenció una mayor heterogeneidad en los modelos producidos, con enfoques que oscilaron entre representaciones informales y adaptaciones de notaciones conocidas como ArchiMate. Esta variabilidad refleja una dificultad estructural que la~\NDITI~busca atender: la ausencia de un lenguaje común que guíe la construcción de modelos de infraestructura en contextos formativos.

En este sentido, los resultados sugieren que la~\NDITI~cumple una función pedagógica concreta: al establecer un vocabulario compartido y una estructura por capas, reduce la dispersión en los enfoques de modelado y facilita la comunicación entre participantes con distintos niveles de experiencia previa. Esta función es especialmente relevante en el contexto del \GRID, donde la diversidad de perfiles y la necesidad de producir representaciones coherentes e interpretables constituyen requisitos habituales del trabajo colaborativo.

La percepción de los participantes refuerza esta lectura. Los comentarios cualitativos recogidos en la encuesta señalan que la notación, con ajustes menores en su representación visual y en la especificidad semántica de algunos elementos, tiene potencial para consolidarse como herramienta de uso regular en actividades de modelado. Uno de los participantes proyectó explícitamente su adopción como estándar, destacando su simplicidad y su capacidad de escalar hacia problemas más complejos. Adicionalmente, el ítem de adecuación para estudiantes en formación obtuvo una de las valoraciones más altas del Eje IV (≈4.25), lo que respalda la pertinencia de la propuesta para el nivel de complejidad propio del contexto educativo en el que fue evaluada.

En conjunto, la evidencia disponible indica que la~\NDITI~representa una contribución viable al ecosistema de herramientas pedagógicas para la enseñanza de Infraestructura de~\TI{}, con un impacto observable en la calidad y consistencia de los modelos producidos por usuarios en etapas tempranas de formación.

\section{Sobre las limitaciones del proyecto}
\noindent
A pesar de los resultados obtenidos, el proyecto presenta limitaciones que deben considerarse al interpretar sus conclusiones y proyectar su alcance.

La más significativa corresponde al tamaño y contexto de la muestra. La validación se realizó con diez subgrupos de tres integrantes cada uno, pertenecientes a un único programa académico y en un único momento de aplicación. Si bien este diseño permitió obtener evidencia comparativa preliminar, restringe la generalización de los resultados a otros contextos institucionales, niveles de formación o perfiles de usuario. En particular, la ausencia de medidas de dispersión estadística formal y el tamaño reducido de la muestra impiden establecer afirmaciones concluyentes sobre la superioridad de la notación más allá del ejercicio específico evaluado.

Una segunda limitación se relaciona con la madurez de la propia notación. Aunque la~\NDITI~fue fundamentada en referentes consolidados como ArchiMate, UML y BPMN, y validada preliminarmente con resultados favorables, algunos aspectos de su especificación requieren refinamiento adicional. Los participantes identificaron ambigüedad semántica en ciertos elementos, dificultades en la diferenciación visual entre capas y una curva de aprendizaje inicial que, aunque no impidió un desempeño adecuado, sugiere que la notación aún no ha alcanzado el nivel de pulimiento necesario para su adopción sin acompañamiento didáctico.

Finalmente, la ausencia de soporte en herramientas de software especializadas representa una limitación práctica relevante. Durante el ejercicio de validación, los participantes debieron construir sus diagramas sin una herramienta que implementara nativamente los elementos y relaciones de la \NDITI, lo que pudo haber introducido variabilidad adicional en los resultados y limita la escalabilidad de su uso en contextos donde la eficiencia y la reutilización de modelos son requisitos importantes.

\section{Sobre los aportes del proyecto}
\noindent
El proyecto genera contribuciones en cuatro dimensiones complementarias que, en conjunto, configuran una propuesta con valor tanto para el ámbito académico inmediato como para líneas de investigación más amplias en el campo del modelado de infraestructuras de TI.

Desde el punto de vista metodológico, el proyecto demuestra la viabilidad de integrar de forma secuencial y coherente tres enfoques complementarios ---caracterización del contexto, mapeo sistemático de literatura y análisis de decisiones--- para el diseño de una notación adaptada a necesidades específicas. Esta articulación metodológica constituye en sí misma un aporte replicable, en tanto ofrece un camino estructurado para abordar problemas similares de diseño de herramientas conceptuales en contextos educativos o institucionales.

En el plano conceptual, la~\NDITI~representa una contribución original al espacio de notaciones de modelado de infraestructura. A diferencia de adoptar un lenguaje existente sin adaptación, la propuesta construye una notación híbrida que toma la estructura multicapa de ArchiMate, la precisión semántica de UML y la orientación comunicativa de BPMN, articulándolas en un sistema coherente orientado a usuarios en formación. Este ejercicio de síntesis conceptual tiene valor más allá del caso específico, al evidenciar que es posible diseñar notaciones especializadas a partir de la combinación razonada de estándares existentes.

El aporte práctico se materializa en los resultados de la validación. El grupo que utilizó la~\NDITI~produjo modelos consistentemente superiores en estructuración, completitud y correspondencia con el enunciado, con una variabilidad interna notablemente menor que la del grupo de referencia. Esto indica que la notación cumple su propósito funcional en condiciones reales de uso, aun en una primera exposición y sin soporte de herramientas especializadas.

Finalmente, el aporte institucional se expresa en la alineación directa de la propuesta con las necesidades identificadas en el programa de Ingeniería de Sistemas y Computación y en el grupo \GRID. La~\NDITI~no es una propuesta genérica trasladada a un contexto específico, sino una herramienta diseñada desde y para ese contexto, lo que incrementa su pertinencia y su potencial de adopción sostenida en procesos formativos y de investigación aplicada.

\section{Conclusión general}
\noindent
El proyecto logró diseñar, fundamentar y validar una propuesta de notación de Infraestructura de~\TI adaptada al contexto académico del programa de Ingeniería de Sistemas y Computación de la Universidad del Quindío. La~\NDITI~responde a necesidades reales identificadas en el proceso de caracterización, se fundamenta en estándares consolidados de modelado y fue evaluada mediante un ejercicio comparativo que permitió obtener evidencia concreta sobre su desempeño y percepción en condiciones de uso real.

Los resultados de la validación indican que la notación contribuye de forma medible a mejorar la calidad de los modelos de infraestructura producidos por usuarios en etapas tempranas de formación. El grupo experimental superó al grupo de referencia en los cuatro criterios evaluados, con diferencias especialmente significativas en estructuración del modelo (4.6 vs. 3.0) y completitud (4.6 vs. 3.4), y con una variabilidad interna notablemente menor, lo que sugiere que la notación establece un piso mínimo de calidad independiente de las capacidades previas de los participantes. La encuesta de percepción complementa estos hallazgos con valoraciones positivas en comprensión, expresividad y utilidad, siendo la facilidad de aprendizaje el único aspecto que señala una oportunidad de mejora concreta.

Estos resultados deben leerse teniendo en cuenta el alcance del estudio: la muestra es reducida, el contexto es específico y la notación se encuentra en una etapa preliminar de desarrollo. Sin embargo, la coherencia entre el desempeño observado y la percepción reportada, junto con la consistencia de los resultados entre subgrupos, ofrece una base de evidencia suficiente para afirmar que la~\NDITI~constituye una propuesta viable y pertinente para el problema que busca resolver.

En términos más amplios, el proyecto demuestra que es posible diseñar herramientas conceptuales adaptadas a contextos educativos específicos a partir de la combinación razonada de estándares existentes, y que su validación en entornos reales de uso genera conocimiento útil tanto para el refinamiento de la herramienta como para la comprensión de las condiciones que facilitan o dificultan su adopción. En este sentido, la~\NDITI~no es únicamente el resultado del proyecto, sino también el punto de partida para una línea de trabajo con potencial de consolidación institucional y académica.